\section{Lecture des normes et anglais technique}

Une part importante du projet a reposé sur la lecture et la compréhension de documents normatifs entièrement rédigés en anglais. Les spécifications \terme{s100} et \terme{s123} ne sont pas de courtes documentations d'API : elles forment un ensemble long, dense et réparti entre la spécification générale, la spécification de produit, les catalogues d'entités, les règles d'encodage, les schémas \terme{xsd} et les fichiers XML associés \cite{iho-s100,iho-s123,iho-s100-registry}. Leur lecture a donc représenté un travail continu d'appropriation, et non une simple consultation ponctuelle.

Il fallait comprendre précisément des notions comme \textit{Feature Catalogue}, \textit{Feature Type}, \textit{Information Type}, \textit{association}, \textit{role}, \textit{multiplicity}, \textit{binding}, \textit{listed value}, \textit{complex attribute} ou \textit{portrayal}. Ces termes devaient ensuite être reliés au modèle de données, aux classes Python générées depuis les schémas et aux concepts métier employés par les spécialistes du Shom. Une mauvaise interprétation d'un terme pouvait conduire à une mauvaise cardinalité, à une association incorrecte ou à une structure de persistance trop éloignée du modèle normatif.

Cette difficulté était renforcée par l'évolution des éditions. La comparaison entre les différentes versions de \terme{s123}, puis l'étude d'autres produits comme S-101, m'ont obligé à identifier les changements de structure, à distinguer les concepts propres à un produit des concepts communs à \terme{s100} et à vérifier leur traduction dans les Feature Catalogues \cite{iho-s101,iho-s123}. Le passage d'une première base directement dérivée de \terme{s123} à un méta-modèle plus générique est en partie issu de cette lecture comparative.

Ce travail constitue un apport concret pour la maîtrise de l'anglais technique. Je n'ai pas seulement consulté des interfaces ou reproduit des exemples de code en anglais : j'ai dû lire des documents normatifs longs, en extraire les règles utiles, reformuler leur contenu en français dans des tableaux de correspondance et des présentations, puis traduire ces règles en choix d'architecture et en code. La documentation de \texttt{xsdata}, FastAPI, PostgreSQL et Pydantic, également majoritairement anglophone, a complété cette pratique lors de l'implémentation \cite{xsdata-docs,fastapi-docs,postgresql-docs,pydantic-docs}.
